\documentclass[11pt, a4paper]{article}

% --- PREAMBOLO ---
\usepackage[a4paper, top=2.5cm, bottom=2.5cm, left=2cm, right=2cm]{geometry}
\usepackage{fontspec}
\usepackage[italian]{babel}

% Impostazione font (Noto Sans per compatibilità e leggibilità)
\babelfont{rm}{Noto Sans}

\usepackage{amsmath}
\usepackage{amssymb}
\usepackage{booktabs}
\usepackage{enumitem}
\usepackage{xcolor}

\title{\textbf{Catalogo Esercizi di Ricerca Operativa}}
\author{Branch and Bound \& Rilassamento Lagrangiano}
\date{Ultimo aggiornamento: Maggio 2026}

\begin{document}

\maketitle

\tableofcontents
\newpage
\section{Catalogo Esercizi: Modellizzazione (Dal Problema Reale al Modello)}
\addcontentsline{toc}{section}{Catalogo Esercizi: Modellizzazione (Dal Problema Reale al Modello)}

% 1. Appello 23 Dicembre 2025 - Versione A
\subsection{Appello: 23 Dicembre 2025 (Versione A)}
Un’azienda alimentare deve produrre $n$ miscele diverse per fare torte. Per produrre le miscele l’azienda usa $m$ materie prime. La produzione di ogni miscela $i$ deve rispettare una ricetta in cui per ogni materia prima $j$ è specificata una quantità $q_{ij}$ da utilizzare per produrre un chilo di miscela. Per ogni miscela $i$, l’azienda ottiene un profitto $p_i$ per ogni chilo prodotto e deve produrre una quantità in chili compresa tra un minimo $a_i$ e un massimo $b_i$. Per ciascuna materia prima $j$, è disponibile la quantità massima $Q_j$. Si vuole definire per ciascuna miscela quanti chili bisogna produrre per massimizzare il profitto complessivo e rispettare i vincoli sopra descritti.
\textit{Scrivere un modello di programmazione lineare.}

% 2. Appello 23 Dicembre 2025 - Versione B
\subsection{Appello: 23 Dicembre 2025 (Versione B)}
Un’azienda elettronica deve produrre $n$ dispositivi diversi utilizzando complessivamente $m$ componenti. Ogni singolo dispositivo $i$ richiede per ogni componente $j$ una specifica quantità $q_{ij}$ e ha un costo di produzione $c_i$. Per ciascun componente $j$, deve essere impiegata una quantità minima $A_j$ e una quantità massima $B_j$. Si vuole definire per ciascun dispositivo quante unità è necessario produrre per minimizzare il costo complessivo e rispettare i vincoli sopra descritti.
\textit{Scrivere un modello di programmazione lineare intera.}

% 3. Appello 23 Dicembre 2025 - Versione C
\subsection{Appello: 23 Dicembre 2025 (Versione C)}
Un’azienda di trasporti deve distribuire un unico prodotto disponibile presso $n$ depositi a $m$ destinazioni diverse. Per ogni deposito $i$ c’è una disponibilità di merce pari a $a_i$ unità, mentre per ogni destinazione c’è una richiesta pari a $b_j$ unità. Ciascun deposito può servire una o più destinazioni, così come ciascuna destinazione può essere servita da uno o più depositi. Il trasporto di una unità di merce dal deposito $i$ alla destinazione $j$ prevede un costo di trasporto pari a $c_{ij}$. Si vuole pianificare come distribuire le merci dai depositi verso le destinazioni minimizzando il costo complessivo rispettando i vincoli sopra descritti.
\textit{Scrivere un modello di programmazione lineare.}

% 4. Appello 23 Dicembre 2025 - Versione D
\subsection{Appello: 23 Dicembre 2025 (Versione D)}
Un’azienda di manutenzione deve assegnare $n$ interventi a $m$ squadre di tecnici. Ogni intervento $i$ richiede $d_i$ giornate e deve essere assegnato a una e una sola squadra; mentre ogni squadra $j$ ha un numero complessivo di giornate a disposizione pari a $Q_j$. Ogni squadra può svolgere un solo intervento alla volta. L’intervento $i$ eseguito dalla squadra $j$ ha un costo complessivo pari a $c_{ij}$. Si vuole pianificare come assegnare gli interventi alle squadre minimizzando il costo complessivo e rispettando i vincoli sopra descritti.
\textit{Scrivere un modello di programmazione lineare (binaria).}

% 5. Appello 13 Gennaio 2026
\subsection{Appello: 13 Gennaio 2026}
Un’azienda deve stabilire come distribuire delle automobili (di un unico modello e un unico allestimento) disponibili presso $n$ depositi verso $m$ concessionari. Per ogni deposito $i$ è indicata una disponibilità di auto pari a $a_i$ unità, mentre per ogni concessionario $j$ è indicata una richiesta pari a $b_j$ unità. Ciascun deposito può servire uno o più concessionari, così come ciascun concessionario può essere servito da uno o più depositi. Il trasporto di una singola auto dal deposito $i$ al concessionario $j$ prevede un costo di trasporto pari a $c_{ij}$, mentre deve essere pagato un costo fisso $f_{ij}$ se dal deposito $i$ al concessionario $j$ è inviata almeno un’auto. Si vuole pianificare come distribuire le auto minimizzando il costo complessivo.
\textit{Scrivere un modello di programmazione lineare (mista-intera).}

% 6. Appello 10 Febbraio 2026
\subsection{Appello: 10 Febbraio 2026}
Un’azienda deve stabilire come distribuire $p$ prodotti disponibili presso $n$ depositi verso $m$ punti di vendita. Per ogni deposito $i$ e per ogni prodotto $k$ è disponibile una quantità $a_{ik}$ unità, mentre ogni punto di vendita $j$ richiede una quantità $b_{jk}$ unità del prodotto $k$. Il trasporto di una unità del prodotto $k$ dal deposito $i$ al punto $j$ prevede un costo di trasporto pari a $c_{ijk}$, mentre deve essere pagato un costo fisso $f_{ij}$ se dal deposito $i$ al punto $j$ è inviata almeno una unità di un qualche prodotto. Si vuole minimizzare il costo complessivo.
\textit{Scrivere un modello di programmazione lineare.}

% 7. Simulazione 1
\subsection{Simulazione 1}
Si consideri un insieme di $n$ oggetti di peso $w_i$, $i = 1, \dots, n$, e un insieme di contenitori di capacità $W$. Si vuole minimizzare il numero di contenitori necessari per contenere tutti gli $n$ oggetti. Ciascun oggetto deve essere caricato tutto in un unico contenitore.
\textit{Scrivere un modello di programmazione lineare intera.}

% 8. Simulazione 2
\subsection{Simulazione 2}
Si consideri un insieme di $m$ risorse disponibili nelle quantità $q_j$, $j = 1, \dots, m$, e un insieme di $n$ prodotti che forniscono un profitto $p_i$, $i = 1, \dots, n$, per ogni unità prodotta. Per ogni unità prodotta ciascun prodotto $i$ consuma $a_{ij}$ unità di ciascuna risorsa $j$. Si vuole decidere quante unità produrre di ciascun prodotto massimizzando il profitto e rispettando le quantità di risorse disponibili. Per ciascun prodotto non è possibile produrre quantità frazionarie.
\textit{Scrivere un modello di programmazione lineare intera.}

% 9. Appello 30 Maggio 2023
\subsection{Appello: 30 Maggio 2023}
Si consideri un insieme di $n$ depositi che hanno una capacità di $Q_i$ unità di merce a settimana e un insieme di $m$ punti vendita che richiedono $q_j$ unità a settimana. Ogni punto vendita deve essere servito da un unico deposito e ogni deposito può servire solo un insieme di punti vendita in cui la somma delle richieste non supera la sua capacità. Per trasportare “tutta” la merce dal deposito $i$ al punto vendita $j$ si paga complessivamente un costo $c_{ij}$. Si vuole minimizzare il costo complessivo.
\textit{Scrivere un modello di programmazione lineare intera.}

% 10. Appello 6 Giugno 2023
\subsection{Appello: 6 Giugno 2023}
Si consideri un insieme di $n$ oggetti di peso $w_i$ quintali e di volume $v_i$ metri cubi. Ad ogni oggetto $i$ è associato un indice di urgenza $u_i$. Gli oggetti devono essere caricati in un insieme di $m$ container che hanno una capacità in peso di $W_j$ quintali and in volume di $V_j$ metri cubi. I container potrebbero non essere sufficienti, per cui si vuole decidere quali oggetti caricare in quali container per massimizzare la somma degli indici di urgenza rispettando i vincoli di capacità.
\textit{Scrivere un modello di programmazione lineare intera.}

% 11. Appello 20 Giugno 2023
\subsection{Appello: 20 Giugno 2023}
Si consideri una rete logistica rappresentata da un grafo $G(V,A)$ in cui $V$ rappresenta punti di produzione, di vendita e di transito. Per ciascun vertice $i \in V$ abbiamo $b_i > 0$ se è un punto di produzione, $b_i < 0$ se è un punto di vendita (richiesta), $b_i = 0$ se è transito. Trasportare una unità di merce sull'arco $(i,j)$ prevede un costo $c_{ij}$ e la quantità non può superare la capacità $u_{ij}$. Si vuole minimizzare il costo complessivo per soddisfare tutte le richieste.
\textit{Scrivere un modello di programmazione lineare.}

% 12. Appello 12 Luglio 2023
\subsection{Appello: 12 Luglio 2023}
Una rete logistica è rappresentata da un grafo $G(V,A)$ bipartito in cui l'insieme dei vertici è partizionato in depositi $V_D$ e punti vendita $V_P$. Per ogni deposito $i \in V_D$ abbiamo una quantità $a_i$ disponibile, per ogni punto vendita $j \in V_P$ una richiesta $b_j$. Ogni arco $(i,j) \in A$ ha un costo $c_{ij}$ per unità trasportata. Si vuole minimizzare il costo complessivo soddisfacendo le richieste con le disponibilità.
\textit{Scrivere un modello di programmazione lineare.}

% 13. Appello 25 Luglio 2023
\subsection{Appello: 25 Luglio 2023}
Un'azienda produce prodotti $P$ utilizzando reparti $R$. Ogni unità del prodotto $i \in P$ genera un profitto $p_i$ e richiede per ciascun reparto $j \in R$ un numero di ore $q_{ij}$. Il numero di unità prodotte per ciascun prodotto $i$ non deve superare la soglia $M_i$. Ciascun reparto $j$ ha a disposizione $Q_j$ ore. Si vuole massimizzare il profitto complessivo.
\textit{Scrivere un modello di programmazione lineare intera.}

% 14. Appello 29 Agosto 2023
\subsection{Appello: 29 Agosto 2023}
Si consideri il problema di caricare prodotti $P$ in contenitori di capacità $W$. I contenitori possono contenere al massimo $K$ prodotti. Ogni prodotto $i \in P$ ha un peso $w_i$ e deve essere caricato interamente in un solo contenitore. Si vuole determinare come caricare tutti i prodotti nel minor numero di contenitori rispettando i vincoli.
\textit{Scrivere un modello di programmazione lineare intera.}

% 15. Appello 14 Settembre 2023
\subsection{Appello: 14 Settembre 2023}
Si consideri il trasporto di una merce in una rete logistica con nodi $N$ (magazzini). In ogni nodo $i$ la merce può essere caricata ($b_i > 0$), scaricata ($b_i < 0$) o transitare ($b_i = 0$). Ogni arco $(i,j)$ ha una capacità $u_{ij}$ e un costo unitario $c_{ij}$. Si vuole minimizzare il costo complessivo rispettando i vincoli di bilancio e capacità.
\textit{Scrivere un modello di programmazione lineare.}

\newpage
\section{Catalogo Esercizi: Simplesso Primale}

% 1. Appello 23 Dicembre 2025 - Versione A
\subsection{Appello: 23 Dicembre 2025 (Versione A - Es. 2)}
Si consideri il seguente problema P:
\[ (P) \begin{cases} \min z = +3x_1 + 2x_2 + x_3 \\ \text{s.t. } +x_1 + x_2 + x_3 \le 6 \\ +2x_1 + x_2 + x_3 \ge 8 \\ x_1, x_2, x_3 \ge 0 \end{cases} \]
Risolvere il problema P utilizzando il metodo del Simplesso Primale.

% 2. Appello 13 Gennaio 2026 - Versione A
\subsection{Appello: 13 Gennaio 2026 (Versione A - Es. 2)}
Si consideri il seguente problema P:
\[ (P) \begin{cases} \min z = -4x_1 + 5x_2 - 2x_3 \\ \text{s.t. } x_1 + x_2 + x_3 = 4 \\ 2x_1 + x_2 + x_3 \le 5 \\ x_1 + x_3 \le 4 \\ x_1, x_2, x_3 \ge 0 \end{cases} \]
Risolvere il problema P utilizzando il metodo del Simplesso Primale.

% 3. Appello 10 Febbraio 2026 - Versione A
\subsection{Appello: 10 Febbraio 2026 (Versione A - Es. 2)}
Si consideri il seguente problema P:
\[ (P) \begin{cases} \min z = +x_1 - 2x_2 + x_3 \\ \text{s.t. } x_1 + x_2 + 3x_3 \le 4 \\ 2x_1 + x_2 + x_3 \le 10 \\ x_1 - x_2 + x_3 = 2 \\ x_1, x_2, x_3 \ge 0 \end{cases} \]
Risolvere il problema P utilizzando il metodo del Simplesso Primale.

% 4. Appello 23 Dicembre 2025 - Versione C
\subsection{Appello: 23 Dicembre 2025 (Versione C - Es. 2)}
Si consideri il seguente problema P:
\[ (P) \begin{cases} \min z = +x_1 + 3x_2 - x_3 \\ \text{s.t. } +x_1 - x_2 + x_3 \le 2 \\ -2x_1 + x_2 + x_3 \ge 3 \\ x_1, x_2, x_3 \ge 0 \end{cases} \]
Risolvere il problema P utilizzando il metodo del Simplesso Primale.

% 5. Simulazione 1
\subsection{Simulazione 1 (Es. 2)}
Si consideri il seguente problema P:
\[ (P) \begin{cases} \min z = -x_1 + 2x_2 + 2x_3 \\ \text{s.t. } +3x_1 + 3x_2 + x_3 \le 2 \\ -x_1 + x_2 + x_3 \ge 1 \\ x_1, x_2, x_3 \ge 0 \end{cases} \]
Risolvere il problema P utilizzando il metodo del Simplesso Primale.

% 6. Appello 30 Maggio 2023
\subsection{Appello: 30 Maggio 2023 (Es. 2)}
Si consideri il seguente problema P:
\[ (P) \begin{cases} \min z = -x_1 + 2x_2 \\ \text{s.t. } 2x_1 - 2x_2 \le 3 \\ x_1 + 2x_2 \ge 6 \\ x_1, x_2 \ge 0 \end{cases} \]
Risolvere il problema P utilizzando il metodo del Simplesso Primale.

% 7. Appello 6 Giugno 2023
\subsection{Appello: 6 Giugno 2023 (Es. 2)}
Si consideri il seguente problema P:
\[ (P) \begin{cases} \min z = -x_1 + 2x_2 + 4x_3 \\ \text{s.t. } x_1 + x_2 + 3x_3 \le 4 \\ x_1 + 2x_2 - x_3 \ge 6 \\ x_1, x_2, x_3 \ge 0 \end{cases} \]
Risolvere il problema P utilizzando il metodo del Simplesso Primale.

% 8. Appello 20 Giugno 2023
\subsection{Appello: 20 Giugno 2023 (Es. 2)}
Si consideri il seguente problema P:
\[ (P) \begin{cases} \min z = +2x_1 - 3x_2 + 3x_3 \\ \text{s.t. } -x_1 + 2x_2 + x_3 \le 2 \\ +4x_1 - 2x_2 - x_3 = 4 \\ x_1, x_2, x_3 \ge 0 \end{cases} \]
Risolvere il problema P utilizzando il metodo del Simplesso Primale.

% 9. Appello 12 Luglio 2023
\subsection{Appello: 12 Luglio 2023 (Es. 2)}
Si consideri il seguente problema P:
\[ (P) \begin{cases} \min z = +x_1 - x_2 - x_3 \\ \text{s.t. } -x_1 + 2x_2 - x_3 \le 2 \\ +2x_1 - 3x_2 + x_3 = 1 \\ x_1, x_2, x_3 \ge 0 \end{cases} \]
Risolvere il problema P utilizzando il metodo del Simplesso Primale.

% 10. Appello 25 Luglio 2023
\subsection{Appello: 25 Luglio 2023 (Es. 2)}
Si consideri il seguente problema P:
\[ (P) \begin{cases} \min z = -4x_1 - x_2 + 3x_3 \\ \text{s.t. } +x_1 + x_2 + 3x_3 \le 6 \\ -2x_1 + x_2 - 4x_3 \ge 3 \\ x_1, x_2, x_3 \ge 0 \end{cases} \]
Risolvere il problema P utilizzando il metodo del Simplesso Primale.

% 11. Appello 29 Agosto 2023
\subsection{Appello: 29 Agosto 2023 (Es. 2)}
Si consideri il seguente problema P:
\[ (P) \begin{cases} \min z = -4x_1 + x_2 + 3x_3 \\ \text{s.t. } -x_1 + x_2 + 3x_3 \le 4 \\ -3x_1 + x_2 - 2x_3 \ge 3 \\ x_1, x_2, x_3 \ge 0 \end{cases} \]
Risolvere il problema P utilizzando il metodo del Simplesso Primale.

% 12. Appello 14 Settembre 2023
\subsection{Appello: 14 Settembre 2023 (Es. 2)}
Si consideri il seguente problema P:
\[ (P) \begin{cases} \min z = +x_1 + 5x_2 + x_3 \\ \text{s.t. } +x_1 + x_2 - 2x_3 \le 4 \\ -x_1 + x_2 + x_3 \ge 2 \\ x_1, x_2, x_3 \ge 0 \end{cases} \]
Risolvere il problema P utilizzando il metodo del Simplesso Primale.

\section{Catalogo Esercizi: Simplesso Duale}

% 1. Appello 23 Dicembre 2025 - Versione B
\subsection{Appello: 23 Dicembre 2025 (Versione B - Es. 2)}
Si consideri il seguente problema P:
\[ (P) \begin{cases} \min z = -5x_1 + x_2 + 4x_3 \\ \text{s.t. } +x_1 + 2x_2 + x_3 \ge 9 \\ +2x_1 + x_2 + 2x_3 \le 12 \\ x_1, x_2, x_3 \ge 0 \end{cases} \]
Risolvere il problema P utilizzando il metodo del Simplesso Duale.

% 2. Appello 23 Dicembre 2025 - Versione D
\subsection{Appello: 23 Dicembre 2025 (Versione D - Es. 2)}
Si consideri il seguente problema P:
\[ (P) \begin{cases} \min z = -x_1 + 3x_2 - 2x_3 \\ \text{s.t. } 4x_1 - 2x_2 - x_3 \ge 4 \\ x_1 + 3x_2 + 2x_3 \le 5 \\ x_1, x_2, x_3 \ge 0 \end{cases} \]
Risolvere il problema P utilizzando il metodo del Simplesso Duale.

% 3. Simulazione 2
\subsection{Simulazione 2 (Es. 2)}
Si consideri il seguente problema P:
\[ (P) \begin{cases} \min z = -x_1 + 2x_2 + 2x_3 \\ \text{s.t. } +3x_1 + 3x_2 + x_3 \le 2 \\ -x_1 + x_2 + x_3 \ge 1 \\ x_1, x_2, x_3 \ge 0 \end{cases} \]
Risolvere il problema P utilizzando il metodo del Simplesso Duale.

% 4. Appello 13 Gennaio 2026 - Es. 3
\subsection{Appello: 13 Gennaio 2026 (Es. 3)}
Si consideri il seguente problema P:
\[ (P) \begin{cases} \min z = +x_1 + 2x_2 - 3x_3 \\ \text{s.t. } x_1 + x_2 + x_3 \le 4 \\ 2x_1 + x_2 + x_3 \le 5 \\ x_1 - x_2 + x_3 = 2 \\ x_1, x_2, x_3 \ge 0 \end{cases} \]
Risolvere il problema P utilizzando il metodo del Simplesso Duale.

% 5. Appello 10 Febbraio 2026 - Es. 3
\subsection{Appello: 10 Febbraio 2026 (Es. 3)}
Si consideri il seguente problema P:
\[ (P) \begin{cases} \min z = x_1 + 2x_2 + 3x_3 \\ \text{s.t. } x_1 + 2x_2 + x_3 = 4 \\ 2x_1 + x_2 + 2x_3 \le 5 \\ x_1 - x_2 + x_3 \le 4 \\ x_1, x_2, x_3 \ge 0 \end{cases} \]
Risolvere il problema P utilizzando il metodo del Simplesso Duale.

% 6. Appello 12 Luglio 2023 - Es. 3
\subsection{Appello: 12 Luglio 2023 (Es. 3)}
Si consideri il seguente problema P:
\[ (P) \begin{cases} \min z = -x_1 + 2x_2 - 2x_3 \\ \text{s.t. } -x_1 + 2x_2 - 2x_3 \ge 4 \\ 2x_1 - x_2 + 2x_3 \le 4 \\ x_1, x_2, x_3 \ge 0 \end{cases} \]
Risolvere il problema P utilizzando il metodo del Simplesso Duale.

% 7. Appello 25 Luglio 2023 - Es. 3
\subsection{Appello: 25 Luglio 2023 (Es. 3)}
Si consideri il seguente problema P:
\[ (P) \begin{cases} \min z = -4x_1 - x_2 + 3x_3 \\ \text{s.t. } x_1 + x_2 + 3x_3 \le 6 \\ -2x_1 + x_2 \ge 3 \\ x_1, x_2, x_3 \ge 0 \end{cases} \]
Risolvere il problema P utilizzando il metodo del Simplesso Duale.

% 8. Appello 14 Settembre 2023 - Es. 3
\subsection{Appello: 14 Settembre 2023 (Es. 3)}
Si consideri il seguente problema P:
\[ (P) \begin{cases} \min z = +3x_1 - x_2 - 2x_3 \\ \text{s.t. } +2x_1 + 2x_2 + x_3 \le 6 \\ +2x_1 + 2x_2 - 2x_3 \ge 3 \\ x_1, x_2, x_3 \ge 0 \end{cases} \]
Risolvere il problema P utilizzando il metodo del Simplesso Duale.

\section{Verifica di Ottimalità (Complementarietà)}

% 1. Appello 23 Dicembre 2025 - Versione A
\subsection{Appello: 23 Dicembre 2025 (Versione A - Es. 3)}
Si consideri il seguente problema P:
\[ (P) \begin{cases} \min z = 5x_1 + x_2 + 4x_3 \\ \text{s.t. } x_1 + 2x_2 + x_3 \le 9 \\ 3x_1 + x_2 + 2x_3 \ge 12 \\ x_1, x_2, x_3 \ge 0 \end{cases} \]
Utilizzando le relazioni di complementarietà verificare se la soluzione $\mathbf{x} = (3,3,0)$ è ottima per il problema P.

% 2. Appello 23 Dicembre 2025 - Versione B
\subsection{Appello: 23 Dicembre 2025 (Versione B - Es. 3)}
Si consideri il seguente problema P:
\[ (P) \begin{cases} \min z = 4x_1 - 2x_2 + x_3 \\ \text{s.t. } x_1 + x_2 + x_3 \ge 6 \\ 2x_1 + x_2 + x_3 \le 8 \\ x_1, x_2, x_3 \ge 0 \end{cases} \]
Utilizzando le relazioni di complementarietà verificare se la soluzione $\mathbf{x} = (2,4,0)$ è ottima per il problema P.

% 3. Appello 23 Dicembre 2025 - Versione C
\subsection{Appello: 23 Dicembre 2025 (Versione C - Es. 3)}
Si consideri il seguente problema P:
\[ (P) \begin{cases} \min z = x_1 - 2x_2 + 3x_3 \\ \text{s.t. } x_1 - x_2 + x_3 \le 5 \\ -2x_1 + 3x_2 + x_3 \le 1 \\ x_1, x_2, x_3 \ge 0 \end{cases} \]
Utilizzando le relazioni di complementarietà verificare se la soluzione $\mathbf{x} = (16,11,0)$ è ottima per il problema P.

% 4. Appello 23 Dicembre 2025 - Versione D
\subsection{Appello: 23 Dicembre 2025 (Versione D - Es. 3)}
Si consideri il seguente problema P:
\[ (P) \begin{cases} \min z = -x_1 - x_2 - x_3 \\ \text{s.t. } x_1 + x_2 + 2x_3 \le 5 \\ -2x_1 + x_2 + 4x_3 \ge 2 \\ x_1, x_2, x_3 \ge 0 \end{cases} \]
Utilizzando le relazioni di complementarietà verificare se la soluzione $\mathbf{x} = (2,0, \frac{3}{2})$ è ottima per il problema P.

% 5. Appello 13 Gennaio 2026
\subsection{Appello: 13 Gennaio 2026 (Es. 4)}
Si consideri il seguente problema P:
\[ (P) \begin{cases} \min z = -x_1 + 5x_2 - 2x_3 \\ \text{s.t. } 4x_1 - 2x_2 + x_3 \le 5 \\ 2x_1 + 3x_2 - x_3 \ge 4 \\ x_1, x_2, x_3 \ge 0 \end{cases} \]
Utilizzando le relazioni di complementarietà verificare se la soluzione $\mathbf{x} = (0,9,23)$ è ottima per il problema P.

% 6. Appello 10 Febbraio 2026
\subsection{Appello: 10 Febbraio 2026 (Es. 4)}
Si consideri il seguente problema P:
\[ (P) \begin{cases} \min z = x_1 - 4x_2 + x_3 \\ \text{s.t. } x_1 - 2x_2 + x_3 \ge 4 \\ -x_1 + 2x_2 + x_3 \le 4 \\ x_1 - x_2 + x_3 \le 6 \\ x_1, x_2, x_3 \ge 0 \end{cases} \]
Utilizzando le relazioni di complementarietà verificare se la soluzione $\mathbf{x} = (8,2,0)$ è ottima per il problema P.

% 7. Simulazione 1
\subsection{Simulazione 1 (Es. 3)}
Si consideri il seguente problema P:
\[ (P) \begin{cases} \min z = -2x_1 - x_2 \\ \text{s.t. } 2x_1 - 3x_2 \ge -4 \\ 3x_1 + 2x_2 \le 7 \\ x_1, x_2 \ge 0 \end{cases} \]
Utilizzando le relazioni di complementarietà verificare se la soluzione $\mathbf{x} = (1,2)$ è ottima per il problema P.

% 8. Simulazione 2
\subsection{Simulazione 2 (Es. 3)}
Si consideri il seguente problema P:
\[ (P) \begin{cases} \min z = -2x_1 + 2x_2 \\ \text{s.t. } -x_1 + 3x_2 \ge 2 \\ 3x_1 + 2x_2 \le 5 \\ x_1, x_2 \ge 0 \end{cases} \]
Utilizzando le relazioni di complementarietà verificare se la soluzione $\mathbf{x} = (1,1)$ è ottima per il problema P.

% 9. Appello 30 Maggio 2023
\subsection{Appello: 30 Maggio 2023 (Es. 3)}
Si consideri il seguente problema P:
\[ (P) \begin{cases} \min z = -x_1 + x_2 + 2x_3 \\ \text{s.t. } 2x_1 - 2x_2 + 3x_3 \le 3 \\ x_1 + 2x_2 + 3x_3 \ge 3 \\ x_1, x_2, x_3 \ge 0 \end{cases} \]
Utilizzando le relazioni di complementarietà verificare se la soluzione $\mathbf{x} = (2, \frac{1}{2}, 0)$ è ottima per il problema P.

% 10. Appello 6 Giugno 2023
\subsection{Appello: 6 Giugno 2023 (Es. 3)}
Si consideri il seguente problema P:
\[ (P) \begin{cases} \min z = -x_1 - x_2 + x_3 \\ \text{s.t. } x_1 - x_2 - x_3 \ge 2 \\ x_1 + 2x_2 - 2x_3 \le 1 \\ x_1, x_2, x_3 \ge 0 \end{cases} \]
Utilizzando le relazioni di complementarietà verificare se la soluzione $\mathbf{x} = (3,0,1)$ è ottima per il problema P.

% 11. Appello 20 Giugno 2023
\subsection{Appello: 20 Giugno 2023 (Es. 3)}
Si consideri il seguente problema P:
\[ (P) \begin{cases} \min z = 2x_1 - 2x_2 - 3x_3 \\ \text{s.t. } x_1 + 2x_2 - 2x_3 \ge -2 \\ x_1 + 3x_2 - x_3 \le 1 \\ x_1, x_2, x_3 \ge 0 \end{cases} \]
Utilizzando le relazioni di complementarietà verificare se la soluzione $\mathbf{x} = (0,1,2)$ è ottima per il problema P.

% 12. Appello 25 Luglio 2023
\subsection{Appello: 25 Luglio 2023 (Es. 3)}
Si consideri il seguente problema P:
\[ (P) \begin{cases} \min z = -4x_1 - x_2 + 3x_3 \\ \text{s.t. } x_1 + x_2 + 3x_3 \le 6 \\ -2x_1 + x_2 \ge 3 \\ x_1, x_2, x_3 \ge 0 \end{cases} \]
Utilizzando le relazioni di complementarietà verificare se la soluzione $\mathbf{x} = (1,5,0)$ è ottima per il problema P.

% 13. Appello 29 Agosto 2023
\subsection{Appello: 29 Agosto 2023 (Es. 3)}
(Testo identico all'Appello del 25 Luglio 2023)
Si consideri il seguente problema P:
\[ (P) \begin{cases} \min z = -4x_1 - x_2 + 3x_3 \\ \text{s.t. } x_1 + x_2 + 3x_3 \le 6 \\ -2x_1 + x_2 \ge 3 \\ x_1, x_2, x_3 \ge 0 \end{cases} \]
Utilizzando le relazioni di complementarietà verificare se la soluzione $\mathbf{x} = (1,5,0)$ è ottima per il problema P.


\section{Parte I: Branch and Bound}
In questa sezione sono presenti gli esercizi in cui viene chiesto esplicitamente di eseguire un'iterazione di Branch and Bound partendo da un tableau ottimo continuo.

\subsection{Appello: 20 Giugno 2023 (Esercizio 4)}
\textbf{Problema P:}
\[ (P) \begin{cases} \min z = -x_1 + 2x_2 - 2x_3 \\ \text{s.t. } -x_1 + 2x_2 - 2x_3 \ge 3 \\ 2x_1 - x_2 + 2x_3 \le 4 \\ x_1, x_2, x_3 \ge 0, \text{ intere} \end{cases} \]

\textbf{Tableau ottimo continuo:}
\[ \begin{array}{c|ccccc|c}
& x_1 & x_2 & x_3 & x_4 & x_5 & RHS \\ \hline
z & 0.000 & 0.000 & 0.000 & -1.000 & 0.000 & 3.000 \\ \hline
x_2 & 1.000 & 1.000 & 0.000 & -1.000 & 1.000 & 7.000 \\
x_3 & 1.500 & 0.000 & 1.000 & -0.500 & 1.000 & 5.500 \\
\end{array} \]

\subsection{Appello: 29 Agosto 2023 (Esercizio 4)}
\textbf{Problema P:}
\[ (P) \begin{cases} \min z = +3x_1 - x_2 - 2x_3 \\ \text{s.t. } 4x_1 + 2x_2 + x_3 \le 6 \\ 3x_1 + x_2 - 2x_3 \ge 2 \\ x_1, x_2, x_3 \ge 0, \text{ intere} \end{cases} \]

\textbf{Tableau ottimo continuo:}
\[ \begin{array}{c|ccccc|c}
& x_1 & x_2 & x_3 & x_4 & x_5 & RHS \\ \hline
z & -4.400 & 0.000 & 0.000 & -0.800 & -0.600 & -3.600 \\ \hline
x_2 & 2.200 & 1.000 & 0.000 & 0.400 & -0.200 & 2.800 \\
x_3 & -0.400 & 0.000 & 1.000 & 0.200 & 0.400 & 0.400 \\
\end{array} \]

\subsection{Schema Meccanico di Risoluzione (Branch and Bound)}
\begin{enumerate}
    \item \textbf{Identificazione Variabile Frazionaria:} Guarda la colonna RHS. Identifica le variabili $x_j$ con valore non intero.
    \item \textbf{Creazione Nodi Figli:} Aggiungi i vincoli $x_j \le \lfloor val \rfloor$ e $x_j \ge \lceil val \rceil$.
    \item \textbf{Ripristino Tableau:} Inserisci il vincolo scelto nel tableau con una nuova slack ($x_6$).
    \item \textbf{Simplesso DUALE:} Usa il simplesso duale per riottimizzare.
\end{enumerate}

\section{Catalogo esercizi sui Tagli di Gomory}

Ecco il catalogo completo degli esercizi sui \textbf{Tagli di Gomory} estratti dalle risoluzioni degli esami passati.

\textbf{Correzione dello Stratega:} Ti sbagli nell'aspettarti 15 esercizi. Come già osservato per il Branch and Bound e il Rilassamento Lagrangiano, l'Esercizio 4 (o talvolta il 3) ruota ciclicamente tra diverse tecniche di Programmazione Lineare Intera. Nelle fonti sono presenti \textbf{4 esercizi d'esame/simulazione} e \textbf{1 esempio guidato} completo nella teoria.

\bigskip

\subsection{Appello: 23 Dicembre 2025 (Versione C - Esercizio 3)}

\textbf{Testo integrale:}

Si consideri il seguente problema $P$:
\[
(P)
\begin{cases}
\min z = x_1 - 2x_2 + 3x_3 \\
x_1 - x_2 + x_3 \le 5 \\
-2x_1 + 3x_2 + x_3 \le 1 \\
x_1, x_2, x_3 \ge 0 \text{ intere}
\end{cases}
\]
Risolvendo il rilassamento continuo del problema $P$ utilizzando il metodo del Simplesso abbiamo ottenuto il seguente tableau ottimo:
\[
\begin{array}{c|ccccc|c}
   & x_1 & x_2 & x_3 & x_4 & x_5 & \text{RHS} \\
\hline
z  & 0.000 & 0.000 & 3.667 & 0.333 & 0.667 & 8.000 \\
\hline
x_1 & 1.000 & 0.000 & 4.000 & 3.000 & 1.000 & 16.000 \\
x_2 & 0.000 & 1.000 & 0.333 & 0.667 & 0.333 & 1.333
\end{array}
\]
a) Definire il Taglio di Gomory relativo alla variabile $x_2$ e dopo averlo aggiunto al tableau riottimizzare.

\bigskip

\subsection{Appello: 23 Dicembre 2025 (Versione D - Esercizio 4)}

\textbf{Testo integrale:}

Si consideri il seguente problema $P$:
\[
(P)
\begin{cases}
\min z = -x_1 - x_2 - x_3 \\
x_1 + x_2 + 2x_3 \le 5 \\
-2x_1 + x_2 + 4x_3 \ge 2 \\
x_1, x_2, x_3 \ge 0 \text{ intere}
\end{cases}
\]
Risolvendo il rilassamento continuo abbiamo ottenuto il seguente tableau ottimo:
\[
\begin{array}{c|ccccc|c}
   & x_1 & x_2 & x_3 & x_4 & x_5 & \text{RHS} \\
\hline
z  & 0.000 & 0.000 & 0.600 & 0.400 & 1.000 & 0.600 \\
\hline
x_2 & -0.600 & 1.000 & 0.200 & -0.200 & -0.000 & 0.200 \\
x_1 & 1.000 & 0.000 & 1.800 & 1.200 & -0.000 & 4.800
\end{array}
\]
a) Definire il Taglio di Gomory relativo alla variabile $x_2$ e dopo averlo aggiunto al tableau riottimizzare.

\bigskip

\subsection{Simulazione 1 (Esercizio 4)}

\textbf{Testo integrale:}

Si consideri il seguente problema $P$:
\[
(P)
\begin{cases}
\min z = -2x_1 - x_2 + 2x_3 \\
x_1 + 2x_2 + 4x_3 \le 4 \\
-2x_1 + x_2 + x_3 \ge 1 \\
x_1, x_2, x_3 \ge 0 \text{ intere}
\end{cases}
\]
Risolvendo il rilassamento continuo abbiamo ottenuto il seguente tableau ottimo:
\[
\begin{array}{c|ccccc|c}
   & x_1 & x_2 & x_3 & x_4 & x_5 & \text{RHS} \\
\hline
z  & -4.000 & 0.000 & 0.000 & -2.000 & -1.000 & -7.000 \\
\hline
x_3 & 1.000 & 0.000 & 1.000 & 0.500 & 0.500 & 2.500 \\
x_2 & -1.000 & 1.000 & 0.000 & -0.250 & 0.250 & 0.250
\end{array}
\]
a) Aggiungere un taglio di Gomory relativo alla variabile $x_3$ (i.e., riga 1 del tableau) e riottimizzare.

\bigskip

\subsection{Appello: 25 Luglio 2023 (Esercizio 4)}

\textbf{Testo integrale:}

Si consideri il seguente problema $P$:
\[
(P)
\begin{cases}
\min z = -2x_1 - x_2 + 2x_3 \\
x_1 + 2x_2 + 4x_3 \le 4 \\
-2x_1 + x_2 + x_3 \ge 1 \\
x_1, x_2, x_3 \ge 0 \text{ intere}
\end{cases}
\]
Il tableau ottimo del rilassamento continuo è:
\[
\begin{array}{c|ccccc|c}
   & x_1 & x_2 & x_3 & x_4 & x_5 & \text{RHS} \\
\hline
z  & 0.000 & 0.000 & -4.600 & -0.800 & -0.600 & -2.600 \\
\hline
x_2 & 0.000 & 1.000 & 1.800 & 0.400 & -0.200 & 1.800 \\
x_1 & 1.000 & 0.000 & 0.400 & 0.200 & 0.400 & 0.400
\end{array}
\]
a) Definire il taglio di Gomory relativo alla variabile $x_2$ e dopo averlo aggiunto al tableau riottimizzare.

\bigskip

\subsection{Schema Meccanico di Risoluzione (Chain of Thought)}

Per non sbagliare la generazione del taglio e la riottimizzazione, imita fedelmente questi passaggi del professore:

\begin{enumerate}
  \item \textbf{Scelta della riga:} Identifica la riga della variabile frazionaria indicata (es. $x_2 = 1.333$ nell'appello di Dicembre).
  \item \textbf{Scomposizione parti frazionarie ($F$):} Per ogni coefficiente $y_{ij}$ e per il termine noto $b_i$ della riga, calcola la parte frazionaria positiva
    \[
      F = \text{val} - \lfloor \text{val} \rfloor.
    \]
    \emph{Attenzione ai negativi:} $\lfloor -0.2 \rfloor = -1 \implies F = -0.2 - (-1) = 0.8$.
  \item \textbf{Costruzione del Taglio:} Il vincolo da aggiungere è sempre
    \[
      -\sum F_{ij} x_j + x_{\text{scarto\_nuova}} = -F_i
    \]
    \emph{Fai attenzione:} Nella sommatoria consideri solo le variabili non in base (quelle che nel tableau originale della riga avevano coefficienti non nulli e non erano 1).
  \item \textbf{Integrazione nel tableau:} Aggiungi il taglio come nuova riga in fondo al tableau e una nuova colonna per la variabile di scarto (es. $x_6$) con coefficiente 1.
  \item \textbf{Riottimizzazione obbligatoria col Simplesso DUALE:}
    \begin{itemize}
      \item Perché? Il nuovo RHS è negativo ($-F_i$), quindi la soluzione è primale non ammissibile.
      \item Poiché partivi da un tableau ottimo, i costi ridotti nella riga 0 sono tutti $\le 0$, soddisfacendo il requisito per il Simplesso Duale.
      \item \textbf{Esecuzione:} Scegli la riga del taglio come riga uscente e applica il rapporto minimo duale per trovare la variabile entrante.
    \end{itemize}
\end{enumerate}

\noindent
\textbf{Nota Operativa:} Il professore richiede spesso di fare un solo pivotaggio col duale per trovare la nuova soluzione ottima o dimostrare che il taglio ha ``separato'' la soluzione frazionaria precedente.

\section{Parte III: Rilassamento Lagrangiano}

\subsection{Appello: 23 Dicembre 2025 (Versione A)}
\textbf{Testo integrale:}
Si consideri il seguente problema P:
\[ (P) \begin{cases} \min z = -4x_1 - 2x_2 + x_3 \\ \text{s.t. } -2x_1 + x_2 + 2x_3 \ge 2 & (1) \\ -x_1 - x_2 + 2x_3 \ge 1 & (2) \\ x_1, x_2, x_3 \in \{0,1\} & (3) \end{cases} \]
a) Rilassare in modo Lagrangiano il vincolo (1) e svolgere la prima iterazione completa del metodo del subgradiente partendo dalle seguenti penalità iniziali:
\begin{itemize}
    \item $i. \lambda = 0$
    \item $ii. \lambda = 2$.
\end{itemize}

\subsection{Appello: 23 Dicembre 2025 (Versione B)}
\textbf{Testo integrale:}
Si consideri il seguente problema P:
\[ (P) \begin{cases} \min z = x_1 - 4x_2 - x_3 \\ \text{s.t. } x_1 - x_2 + 2x_3 \ge 1 & (1) \\ 2x_1 - x_2 - x_3 \ge 1 & (2) \\ x_1, x_2, x_3 \in \{0,1\} & (3) \end{cases} \]
a) Rilassare in modo Lagrangiano il vincolo (1) e svolgere la prima iterazione completa del metodo del subgradiente partendo dalle seguenti penalità iniziali:
\begin{itemize}
    \item $i. \lambda = 0$
    \item $ii. \lambda = 2$.
\end{itemize}

\subsection{Simulazione 2}
\textbf{Testo integrale:}
Si consideri il seguente problema P:
\[ (P) \begin{cases} \min z = x_1 + x_2 + 3x_3 \\ \text{s.t. } x_1 - x_2 - x_3 \le 1 & (1) \\ x_1 - x_2 + x_3 \ge 1 & (2) \\ x_1, x_2, x_3 \in \{0,1\} & (3) \end{cases} \]
a) Rilassare in modo Lagrangiano il vincolo (1) e svolgere la prima iterazione completa del metodo del subgradiente partendo dalle seguenti penalità iniziali:
\begin{itemize}
    \item $i. \lambda_1 = 0$
    \item $ii. \lambda_1 = -2$.
\end{itemize}

\subsection{Appello: 30 Maggio 2023}
\textbf{Testo integrale:}
Si consideri il seguente problema P:
\[ (P) \begin{cases} \min z = 2x_1 - x_2 + 3x_3 \\ \text{s.t. } x_1 + x_2 + x_3 \ge 2 & (1) \\ x_1 - x_2 + x_3 \ge 1 & (2) \\ x_1, x_2, x_3 \in \{0,1\} & (3) \end{cases} \]
a) Rilassare in modo Lagrangiano il vincolo (1) e svolgere la prima iterazione completa del metodo del subgradiente partendo dalle seguenti penalità iniziali:
\begin{itemize}
    \item $i. \lambda_1 = 0$
    \item $ii. \lambda_1 = 2$.
\end{itemize}

\subsection{Appello: 12 Luglio 2023}
\textbf{Testo integrale:}
Si consideri il seguente problema P:
\[ (P) \begin{cases} \min z = 2x_1 + 4x_2 + x_3 \\ \text{s.t. } x_1 + x_2 + x_3 \ge 2 & (1) \\ 2x_1 - x_2 + x_3 \le 1 & (2) \\ x_1, x_2, x_3 \in \{0,1\} & (3) \end{cases} \]
a) Rilassare in modo Lagrangiano il vincolo (1) e svolgere la prima iterazione completa del metodo del subgradiente partendo dalle seguenti penalità iniziali:
\begin{itemize}
    \item $i. \lambda_1 = 0$
    \item $ii. \lambda_1 = 4$.
\end{itemize}

\subsection{Appello: 14 Settembre 2023}
\textbf{Testo integrale:}
Si consideri il seguente problema P:
\[ (P) \begin{cases} \min z = 2x_1 - 2x_2 + x_3 \\ \text{s.t. } -x_1 + x_2 + 2x_3 \ge 2 & (1) \\ 2x_1 - x_2 + x_3 \ge 1 & (2) \\ x_1, x_2, x_3 \in \{0,1\} & (3) \end{cases} \]
a) Rilassare in modo Lagrangiano il vincolo (1) e svolgere la prima iterazione completa del metodo del subgradiente partendo dalle seguenti penalità iniziali:
\begin{itemize}
    \item $i. \lambda_1 = 0$
    \item $ii. \lambda_1 = 3$.
\end{itemize}

\subsection{Schema Meccanico di Risoluzione (Chain of Thought)}
\begin{enumerate}
    \item \textbf{Segno di $\lambda$:} Se rilassi un vincolo $\ge$, allora $\lambda \ge 0$. Se rilassi un vincolo $\le$, allora $\lambda \le 0$.
    \item \textbf{Scrittura $z_{LR}(\lambda)$:} Combina la funzione obiettivo originale con il vincolo rilassato nella forma: $\min z + \lambda(b - Ax)$.
    \item \textbf{Risoluzione del Sottoproblema:} Con $\lambda$ fissato, minimizza la nuova funzione rispetto ai vincoli rimanenti (solitamente la binarietà $\{0,1\}$ e un altro vincolo facile).
    \item \textbf{Calcolo del Subgradiente ($s$):} Usa la formula $s = b - Ax$.
    \item \textbf{Aggiornamento:} La nuova penalità sarà $\lambda_{new} = \max\{0, \lambda + \theta s\}$ (se $\lambda \ge 0$) o $\lambda_{new} = \min\{0, \lambda + \theta s\}$ (se $\lambda \le 0$).
    \item \textbf{Verifica Ottimalità:} Se la soluzione $x$ trovata è ammissibile per il problema originale e $\lambda \cdot s = 0$, allora hai trovato l'ottimo.
\end{enumerate}

\end{document}
